\chapter{实数轴上的紧致性与连通性}

\section*{A组　覆盖、紧致与极值}
\addcontentsline{toc}{section}{A组　覆盖、紧致与极值}

\begin{solution}{18.1}
对 \((0,1)\) 取 \(U_n=(1/n,1)\)；对 \([0,\infty)\) 取 \(U_n=(-1,n)\)；对 \(\R\) 取 \(U_n=(-n,n)\)。三族都分别覆盖目标集。任意有限子族只需取最大指标 \(N\)，其并分别仍漏掉靠近 \(0\) 的正点、大于等于 \(N\) 的点或绝对值不小于 \(N\) 的点。
\end{solution}

\begin{solution}{18.2}
令 \(E\) 为所有使 \([a,x]\) 可有限覆盖的 \(x\)。由含 \(a\) 的开集可知 \(E\ne\varnothing\)，且 \(E\subset[a,b]\)，故 \(c=\sup E\) 存在。若 \(c<b\)，取覆盖成员 \(V\ni c\) 及 \((c-\rho,c+\rho)\subset V\)。从 \(E\) 中取 \(x>c-\rho/2\)，覆盖 \([a,x]\) 的有限成员再加 \(V\) 会覆盖到 \(c+\rho/2\)，与 \(c\) 为上界矛盾。因此 \(c=b\)。最后取含 \(b\) 的开集和一个足够靠近 \(b\) 的 \(E\) 中点，得到 \([a,b]\) 的有限覆盖。
\end{solution}

\begin{solution}{18.3}
上确界不必属于原集合。例如 \(\sup(0,1)=1\notin(0,1)\)。因此 \(c=b\) 只说明可有限覆盖的左端区间能任意接近 \(b\)，尚未说明整个 \([a,b]\) 可覆盖。利用含 \(b\) 的开集向左覆盖一小段，再与某个足够靠近 \(b\) 的已覆盖区间拼接，才得到 \(b\in E\)。
\end{solution}

\begin{solution}{18.4}
若 \(K\) 闭且有界，取 \(K\subset[-M,M]\)。给定 \(K\) 的开覆盖，添入开集 \(\R\setminus K\)，便得到 \([-M,M]\) 的开覆盖。闭区间定理给出有限子覆盖。删去 \(\R\setminus K\) 后，余下原覆盖成员仍覆盖 \(K\)。
\end{solution}

\begin{solution}{18.5}
若 \(K\) 无界，则 \(\{(-n,n):n\in\N\}\) 是 \(K\) 的开覆盖。任意有限子族的并为某个 \((-N,N)\)，不能覆盖无界的 \(K\)，与有限覆盖性质矛盾。
\end{solution}

\begin{solution}{18.6}
若 \(K\) 不闭，存在 \(a\in\overline K\setminus K\)。令
\[
 U_n=\{x:\abs{x-a}>1/n\}.
\]
因 \(a\notin K\)，这些开集覆盖 \(K\)。有限子族的并等于其中指标最大的一个 \(U_N\)，但 \(a\) 是 \(K\) 的聚点，存在 \(x\in K\) 满足 \(\abs{x-a}<1/N\)，故 \(x\notin U_N\)。矛盾。
\end{solution}

\begin{solution}{18.7}
覆盖法：每点连续性给出邻域 \(U_x\)，使其中函数值与 \(f(x)\) 相差小于 \(1\)。取有限子覆盖 \(U_{x_1},\ldots,U_{x_m}\)，以
\[
 1+\max_j\abs{f(x_j)}
\]
为全局界。

序列法：若无界，选 \(x_n\) 使 \(\abs{f(x_n)}>n\)。闭区间中的点列有收敛子列，其极限仍在闭区间；连续性使像子列收敛而有界，与构造矛盾。
\end{solution}

\begin{solution}{18.8}
由有界性令 \(M=\sup f([a,b])\)。选 \(x_n\) 使
\[
 M-1/n<f(x_n)\le M.
\]
闭区间序列紧致性给出 \(x_{n_k}\to x_*\in[a,b]\)。连续性与夹逼给出
\[
 f(x_*)=\lim_k f(x_{n_k})=M.
\]
最大值取得。对 \(-f\) 重复，得到最小值。
\end{solution}

\begin{solution}{18.9}
删去闭性：\(f(x)=x\) 在 \((0,1)\) 上不取极值。删去有界性：\(f(x)=x\) 在闭集 \(\R\) 上无界。删去连续性：在 \([0,1]\) 上令
\[
 f(0)=0,\qquad f(x)=1/x\quad(x>0),
\]
则无界，因而没有最大值。
\end{solution}

\section*{B组　连通与介值}
\addcontentsline{toc}{section}{B组　连通与介值}

\begin{solution}{18.10}
设区间 \(I=A\cup B\) 被两个非空相对开集分离，取 \(a\in A,b\in B\)，不妨 \(a<b\)。令 \(c=\sup(A\cap[a,b])\)。区间性质保证 \(c\in I\)。若 \(c\in A\)，相对开性使 \(c\) 右侧仍有 \(A\) 中点，违背上界；若 \(c\in B\)，相对开性在 \(c\) 左侧排除 \(A\) 中点，违背上确界逼近。故不存在分离。
\end{solution}

\begin{solution}{18.11}
若连通集 \(E\) 不具区间性质，存在 \(x<z<y\)，其中 \(x,y\in E\)、\(z\notin E\)。则
\[
 E\cap(-\infty,z),\qquad E\cap(z,\infty)
\]
是两个非空、不交的相对开集，其并为 \(E\)，构成分离，矛盾。
\end{solution}

\begin{solution}{18.12}
取无理数 \(\alpha\)。则
\[
 \Q\cap(-\infty,\alpha),\qquad \Q\cap(\alpha,\infty)
\]
是 \(\Q\) 的非空相对开分离。对无理数集，取任意有理数 \(r\)，用 \(r\) 左右两个半轴与无理数集相交即可分离。两集合虽都在实线中稠密，仍不连通。
\end{solution}

\begin{solution}{18.13}
设 \(f(a)<0<f(b)\)，令 \(E=\{x:f(x)<0\}\)、\(c=\sup E\)。若 \(f(c)<0\)，连续性保证 \(c\) 右侧仍有负值，产生大于 \(c\) 的 \(E\) 中点；若 \(f(c)>0\)，连续性保证 \(c\) 左侧邻域为正，但上确界性质要求 \(E\) 中点从左任意逼近 \(c\)。两次连续性分别用来保持负号与保持正号，故只能 \(f(c)=0\)。
\end{solution}

\begin{solution}{18.14}
给定介于 \(f(x_0),f(x_1)\) 之间的 \(y\)，令 \(g(x)=f(x)-y\)。若 \(y\) 等于某端点值，直接取该端点；否则 \(g\) 在两端严格异号。对 \(g\) 应用零点定理得到 \(g(c)=0\)，即 \(f(c)=y\)。
\end{solution}

\begin{solution}{18.15}
任取 \(u,v\in f(I)\)，选 \(x_0,x_1\in I\) 使 \(f(x_0)=u,f(x_1)=v\)。对每个介于 \(u,v\) 之间的 \(y\)，介值定理给出位于 \(x_0,x_1\) 之间的 \(c\in I\)，使 \(f(c)=y\)。所以 \(f(I)\) 具有区间性质。
\end{solution}

\begin{solution}{18.16}
若存在 \(x_0,x_1\) 使整数值 \(f(x_0)<f(x_1)\)，则两整数之间可取非整数 \(y\)。介值定理会给出 \(c\) 使 \(f(c)=y\)，与值域包含于 \(\Z\) 矛盾。因此任意两点函数值相同，函数为常值。
\end{solution}

\begin{solution}{18.17}
令 \(g(x)=f(x)-x\)。函数 \(g\) 连续，且
\[
 g(a)=f(a)-a\ge0,\qquad
 g(b)=f(b)-b\le0.
\]
若某端点值为零即完成；否则两端异号，由零点定理存在 \(c\) 使 \(g(c)=0\)，即 \(f(c)=c\)。
\end{solution}

\begin{solution}{18.18}
设多项式次数为奇数，首项为 \(a_{2m+1}x^{2m+1}\)。低次项与首项之比在 \(\abs x\to\infty\) 时趋零，因此当 \(\abs x\) 足够大时，多项式符号与首项相同。正负无穷方向上首项符号相反，故可取 \(A>0\) 使 \(p(-A)p(A)<0\)。由连续性和零点定理，\(p\) 在 \((-A,A)\) 内有实根。
\end{solution}

\section*{C组　结构、边界与项目}
\addcontentsline{toc}{section}{C组　结构、边界与项目}

\begin{solution}{18.19}
给定 \(f(K)\) 的开覆盖 \(\{V_\lambda\}\)。连续性使 \(f^{-1}(V_\lambda)\) 在 \(K\) 中相对开，并覆盖 \(K\)。写成 \(K\cap U_\lambda\) 后，由 \(K\) 的有限覆盖性质取有限多个原像覆盖 \(K\)。对应的有限多个 \(V_\lambda\) 覆盖 \(f(K)\)。
\end{solution}

\begin{solution}{18.20}
若 \(f(E)=A\cup B\) 可由两个非空、不交的相对开集分离，则
\[
 f^{-1}(A),\qquad f^{-1}(B)
\]
是 \(E\) 的两个非空、不交、相对开的集合，且并为 \(E\)。这与 \(E\) 连通矛盾，故连续像连通。
\end{solution}

\begin{solution}{18.21}
设 \(S\subset[a,b]\) 无限。把区间二分，至少一个半区间含 \(S\) 的无限多个点；递归选择这种闭半区间，得到套区间 \(I_n\)，长度趋零且每个都含 \(S\) 的无限点。套区间定理给出唯一公共点 \(c\)。每个 \(c\) 的邻域包含某个 \(I_n\)，其中有异于 \(c\) 的 \(S\) 中点，故 \(c\) 是聚点。

若闭区间有一个无有限子覆盖的开覆盖，可递归选择不能有限覆盖的半区间。其公共点属于某覆盖成员，该开集会包含充分小的套区间，与该半区间不能有限覆盖矛盾。
\end{solution}

\begin{solution}{18.22}
设 \(f:K\to Y\subset\R\) 是连续双射，\(K\) 紧致。要证 \(f^{-1}\) 连续，只需证它把 \(K\) 中闭集的像变成 \(Y\) 中闭集。若 \(C\subset K\) 闭，则 \(C\) 也是紧集；连续像 \(f(C)\) 紧，因而在实线中闭。于是 \(f\) 是闭映射，其逆映射连续。
\end{solution}

\begin{solution}{18.23}
证明链可组织为：
\[
 \text{闭且有界}\Rightarrow\text{序列紧致}
\]
使用 Bolzano--Weierstrass 与闭性；
\[
 \text{序列紧致}\Rightarrow\text{有限覆盖性质}
\]
反设无有限子覆盖，从逐渐增大的有限子族外选点，再由收敛子列落入某覆盖成员导致矛盾；
\[
 \text{有限覆盖性质}\Rightarrow\text{闭且有界}
\]
分别使用覆盖 \((-n,n)\) 和避开缺失聚点的覆盖。实数完备性进入 Bolzano--Weierstrass 或闭区间有限覆盖定理的证明。
\end{solution}

\begin{solution}{18.24}
对每个 \(x\in K\)，设连续性在精度 \(\varepsilon\) 下给出半径 \(\delta_x\)。以半径 \(\delta_x/2\) 的邻域覆盖 \(K\)，取有限子覆盖。令 \(\lambda\) 为这些半径的一半中的最小值。任何直径小于 \(\lambda\) 且与 \(K\) 相交的局部片段会落入某个原邻域，从而其函数振荡小于所需精度。这种统一的 \(\lambda\) 是 Lebesgue 数思想，也是 Heine--Cantor 定理的覆盖证明基础。
\end{solution}

\begin{solution}{18.25}
任取 \(x_n\in K_n\)。因 \(K_1\) 紧，存在子列 \(x_{n_k}\to x\in K_1\)。固定 \(m\)。当 \(n_k\ge m\) 时，由递减性 \(x_{n_k}\in K_{n_k}\subset K_m\)。集合 \(K_m\) 闭，故极限 \(x\in K_m\)。这对每个 \(m\) 成立，所以
\[
 x\in\bigcap_{m=1}^{\infty}K_m.
\]
\end{solution}

\begin{solution}{18.26}
零点集
\[
 Z=f^{-1}(\{0\})
\]
在 \(K\) 中闭，因此是紧集。若 \(f\) 处处非零，则连续函数 \(\abs f\) 在紧集 \(K\) 上取得最小值 \(m\ge0\)。若 \(m=0\)，存在 \(x_0\in K\) 使 \(\abs{f(x_0)}=0\)，与处处非零矛盾，故 \(m>0\)。
\end{solution}
